\section{Research Gap and Study Rationale}


The literature establishes three layers of context.
First, \gls{lusc} is a molecularly heterogeneous disease with reproducible transcriptomic subtypes and a narrow actionable-driver landscape, creating a compelling need for improved molecular characterisation \parencite{Wilkerson2010, Hammerman2012, Lau2022-hy}. Second, a large and growing body of work uses machine learning to classify cancer subtypes and identify biomarkers, but the vast majority of studies use a single omics platform and a single model architecture, limiting the biological and methodological conclusions that can be drawn \parencite{Picard2021, Nicora2020}. Third, recent advances in deep learning and foundation models have expanded what is possible with large-scale pre-trained models, but the data requirements and interpretability constraints of these approaches make them poorly suited to the high-dimensional, low-sample-size setting of TCGA modality-specific cohorts.
 
Three specific gaps motivate this thesis:
 
\textbf{Gap 1: No systematic modality-by-modality comparison for established \gls{lusc} subtypes.} Most existing classification studies either use expression data alone, integrate modalities before modelling, or evaluate modalities in separate publications with different preprocessing and evaluation frameworks. It is therefore unclear how faithfully methylation, \gls{mirna}, CNV, and RPPA independently recover the Wilkerson expression-defined subtype taxonomy, and whether their discriminative performance reflects genuine biological signal or platform-specific artefacts.
 
\textbf{Gap 2: Architecture-dependent feature selection is uncharacterised in multi-omics \gls{lusc} classification.} Linear models and tree-based models are known to implement different inductive biases, but the consequences for molecular feature selection in cancer omics have not been systematically quantified. If biomarker nominations are model-architecture-dependent- as is hypothesised- then single-model interpretation studies present an incomplete and potentially misleading view of the discriminative molecular landscape.
 
\textbf{Gap 3: SHAP fold stability across modalities in a small-sample multi-class cancer setting is uncharacterised.} Feature importance stability has been studied in general high-dimensional settings \parencite{Kalousis2007}, but modality-specific and architecture-specific stability patterns in the TCGA-\gls{lusc} classification regime have not been documented. Understanding which modalities produce more reproducible feature rankings is necessary for calibrating confidence in biomarker nominations.
 
This thesis addresses all three gaps through a unified computational framework. Five classifiers are trained and evaluated independently on five TCGA-\gls{lusc} omics modalities, with SHAP-based feature interpretation, fold-stability analysis, cross-model Jaccard similarity across modality-classifier combinations. The result is a modality-resolved and architecture-resolved account of \gls{lusc} subtype biology, which establishes which molecular layers are most informative, which model families converge on shared biological signals, and what interpretive conclusions can be drawn with confidence from small-sample multi-omics classification.
